\subsection{Mô hình OLS Rút Gọn (OLS Chọn Biến)}

Hàm \texttt{run\_ols\_selected} có nhiệm vụ xây dựng một mô hình OLS tinh gọn bằng cách chủ động loại bỏ thêm các biến rác để tìm ra một tập hợp các đặc trưng tinh túy nhất. Trái ngược với mô hình OLS cơ bản, quá trình này được thực hiện qua hai màng lọc cực kỳ khắt khe:

\subsubsection{Màng lọc 1: Lọc Đa cộng tuyến ngặt nghèo (Ngưỡng VIF $\le 5.0$)}
Mặc dù ở bước tiền xử lý (\texttt{data\_pipeline.py}) đã áp dụng lọc VIF với ngưỡng 10.0, ở mô hình này, thuật toán \texttt{run\_vif\_feature\_selection} tiếp tục hạ ngưỡng VIF xuống \textbf{5.0}. Đây là tiêu chuẩn ngặt nghèo hơn rất nhiều, nhằm bảo đảm các biến độc lập được đưa vào mô hình gần như không có sự chồng chéo thông tin hay tương quan tuyến tính nào với nhau.

\subsubsection{Màng lọc 2: Lọc theo ý nghĩa thống kê (p-value $\le 0.05$)}
Sau khi qua màng lọc VIF, dữ liệu được đưa vào chạy hồi quy OLS lần thứ nhất. Hàm \texttt{coef\_inference} được sử dụng để tính toán giá trị \textbf{p-value} cho từng hệ số hồi quy ($\beta$).
\begin{itemize}
    \item \textbf{Ý nghĩa của p-value:} Thể hiện xác suất mà một biến thực chất không có tác động gì đến biến mục tiêu, nhưng do ngẫu nhiên hệ thống lại tính toán ra một hệ số khác 0.
    \item \textbf{Tiêu chí lọc:} Thuật toán lặp qua từng biến, bất kỳ biến nào có $p\text{-value} > 0.05$ (tức là không có ý nghĩa thống kê ở mức ý nghĩa 5\%) sẽ bị loại bỏ. Chỉ những biến thỏa mãn $p\text{-value} \le 0.05$ (cùng với biến hằng số - Intercept) mới được giữ lại.
\end{itemize}

\subsubsection{Huấn luyện mô hình tinh hoa}
Sau khi trải qua hai màng lọc trên, danh sách biến ban đầu được rút gọn lại thành tập hợp các đặc trưng thực sự quan trọng. Hệ thống sẽ thực hiện huấn luyện mô hình OLS một lần cuối cùng chỉ trên các biến tinh hoa này. Mô hình cuối cùng này sau đó được sử dụng để dự đoán trên tập kiểm thử (Test set) và tính toán các chỉ số đánh giá như MAE, RMSE, và $R^2$.

\textbf{Kết luận:} Phương pháp OLS Chọn Biến ứng dụng kỹ thuật loại bỏ đa cộng tuyến khắt khe kết hợp với kiểm định giả thuyết (Backward Elimination). Kết quả mang lại là một mô hình tinh gọn, minh bạch, tập trung vào các yếu tố cốt lõi có tác động thực sự đến lương của cầu thủ.
